\documentclass[11pt]{article}

\usepackage{hyperref}
\usepackage{xcolor}
\usepackage{calc}
\usepackage{graphicx}
\usepackage{tikz}
\usepackage{fontspec}
\usepackage{fontawesome5}
\usepackage{titlesec}
\usepackage{enumitem}
\usepackage{fancybox}
\usepackage{tabularx}
\usepackage{adjustbox}
\hypersetup{hidelinks}

%%%%%%%%%%%%%%%%%%%%
% 设置
%%%%%%%%%%%%%%%%%%%%

\setlength{\parindent}{0pt}					% 取消全局段落缩进
\pagenumbering{gobble}						% 取消页码显示
\setlist[itemize]{nosep                     % 取消 itemize 的默认间距
    , before={\vspace*{-\parskip}}          % 取消 itemize 和后续段落之间的空白
    , leftmargin=*}		                    % 取消 itemize 的左边距
\setlist[enumerate]{leftmargin=*}	        % 取消 enumerate 的左边距
\renewcommand{\arraystretch}{1.2}           % 设置表格行间距
\linespread{1.25}                           % 设置正文行间距

\titleformat{\section}					    % 将原标题前面的数字取消了
  {\LARGE\bfseries\raggedright} 		      % 字体改为 LARGE，bold，左对齐
  {}{0em}                      			  % 可用于添加全局标题前缀
  {}                           			  % 可用于添加代码
  [{\color{secondary_color}\titlerule}]     % 标题下方加一条线
\titlespacing*{\section}{0cm}{*1.2}{*1.2}	% 标题左边留白，上方，下方

\usepackage[
	a4paper,
	left=1.2cm,
	right=1.2cm,
	top=1.5cm,
	bottom=1cm,
	nohead
]{geometry}                                 % 页面边距设置

% 字体设置
\setmainfont[
    Path=fonts/,
    Extension=.otf,
    BoldFont=*-Bold,
]{NotoSerifSC}

% 自定义颜色
\definecolor{primary_color}{RGB}{0, 91, 172}        % 北航蓝
\definecolor{secondary_color}{RGB}{125, 123, 122}   % 银河灰

\newlength{\iconwidth}
\setlength{\iconwidth}{1.5em}                   % 设置 section 标题部分图标占用的宽度

%%%%%%%%%%%%%%%%%%%%
% 文章内容
%%%%%%%%%%%%%%%%%%%%

% 学院
\newcommand{\school}{自动化科学学院 | School of Automation Science} 

% 联系方式
\newcommand{\contact}{
    % 根据个人喜好选择字号
    % \small                % 小
    \footnotesize           % 更小
    % \scriptsize           % 再小一号
    \textcolor{white}{
        % 邮箱
        \faEnvelope \quad \href{mailto:youremail@buaa.edu.cn}{name@buaa.edu.com}
        \hspace{4em}
        % 手机号
        \faPhone \quad  1xx-xxxx-xxxx
        % % 微信
        % \hspace{4em}
        % \faWeixin \quad wangyunpeng_1919810
        % % QQ
        % \hspace{4em}
        % \faQq \quad 2xxxxxxxxx
        % GitHub
        \hspace{4em}
        \faGithub \quad \href{https://github.com/v0yager33/resume-latex-buaa}{GitHub 地址}
    }
}

\begin{document}

    %%%%%%%%%%%%%%%%%%%%
    % 页眉、页脚和背景（如果有多页简历，请把页眉页脚和背景复制粘贴到第二页的内容之前）
    %%%%%%%%%%%%%%%%%%%%

    % 页眉：校标组合+学院名
    \begin{tikzpicture}[remember picture, overlay]
        \node[anchor=north, inner sep=0pt](header) at (current page.north){
            \includegraphics[width=\paperwidth]{images/header.png}
        };
        \node[anchor=west](school_logo) at (header.west){
            \hspace{0.5cm}
            \includegraphics[width=0.25\textwidth]{images/banner.png}
        };
        \node[anchor=east](school_name) at(header.east){
            \textcolor{white}{\textbf{\school}}
            \hspace{0.5cm}
        };
    \end{tikzpicture}
    \vspace{-3.5em}

    % 页脚，联系方式
    \begin{tikzpicture}[remember picture, overlay]
        \node[anchor=south, inner sep=0pt](footer) at (current page.south){
            \includegraphics[width=\paperwidth]{images/footer.png}
        };
        % 联系方式
        \node[anchor=center] at(footer.center){\contact};
    \end{tikzpicture}

    % 背景
    \begin{tikzpicture}[remember picture, overlay]
        \node[opacity=0.05] at(current page.center){
            \includegraphics[width=0.7\paperwidth, keepaspectratio]{images/logo.png}
        };
    \end{tikzpicture}

    %%%%%%%%%%%%%%%%%%%%
    % 简历正文
    %%%%%%%%%%%%%%%%%%%%

    \begin{minipage}[t]{0.75\textwidth}
        % 个人信息
        % \faGraduationCap这类\fa开头的都是font awesome里的logo，想换成其他logo的话，可以看一下附带的fontawsome.pdf，自行替换。
        \begin{minipage}[t]{\textwidth}
        \section[个人信息]{\makebox[\iconwidth][c]{\color{primary_color}{\faAddressCard}}\quad 个人信息}
        \begin{minipage}[t]{0.56\textwidth}
            \textbf{姓\qquad 名}：你的名字
            
            \vspace{0.5em}
            \textbf{出生年月}：2026 年 11 月

            \vspace{0.5em}
            \textbf{联系方式}：1XX-XXXX-XXXX

            \vspace{0.5em}
            \textbf{电子邮箱}：name@mail.buaa.com
        \end{minipage}
        \begin{minipage}[t]{0.35\textwidth}
            \textbf{性\qquad 别}：男

            \vspace{0.5em}
            \textbf{民\qquad 族}：汉
            
            \vspace{0.5em}
            \textbf{政治面貌}：共青团员

            \vspace{0.5em}
            \textbf{籍\qquad 贯}：江苏省南京市
        \end{minipage}
        \vspace{1.2em}
        \end{minipage}
    \end{minipage}
    \hfill
    % 右半边，照片，比例占行宽20%
    \begin{minipage}[t]{0.2\textwidth}
        \vspace{0.2em} % 照片上侧内容
        \setlength{\fboxsep}{0pt}
        \doublebox{\includegraphics[width=0.8\linewidth]{images/avatar.png}}
    \end{minipage}
    \begin{minipage}[t]{\textwidth}
        % 教育背景
        \begin{minipage}[t]{\textwidth}
        \section[教育背景]{\makebox[\iconwidth][c]{\color{primary_color}{\faGraduationCap}}\quad 教育背景}
                
        \vspace{0.5em}
        {\large \textbf{北京航空航天大学}}，博士 \hfill 2024年9月--2027年7月
        \begin{itemize}
            \item \textbf{专\qquad 业}：模式识别与智能系统（081102）    \hfill 师从\ 名字\ 教授
            \item \textbf{研究方向}：计算机视觉、目标检测、多模态交互、具身智能、三维重建\ 等。
            \item \textbf{主修课程}：机器学习、算法设计与分析、虚拟现实技术、工业互联网、最优化理论与方法\ 等。
        \end{itemize}
        
        \vspace{0.5em}
        {\large \textbf{北京航空航天大学}}，大专 \hfill 2020年9月--2024年6月
        
        % 双栏区域
        \vspace{-0.8em} 
        \noindent\begin{minipage}[t]{0.42\linewidth}
            \begin{itemize}
                \setlength{\leftmargin}{1.em} % 调整缩进以匹配下方列表
                \setlength{\itemsep}{0pt}      % 去除项间距
                \setlength{\topsep}{0pt}       % 去除顶部间距
                \item \textbf{专\qquad 业}：虚拟现实技术
            \end{itemize}
        \end{minipage}%
        \hfill % 中间填充空隙
        \begin{minipage}[t]{0.48\linewidth}
            \begin{itemize}
                \setlength{\leftmargin}{1.2em}
                \setlength{\itemsep}{0pt}
                \setlength{\topsep}{0pt}
                \item \textbf{绩点排名}：3.8 / 4\  (1 / 514)，国家奖学金
            \end{itemize}
        \end{minipage}
        \vspace{0.5em}
        
        \begin{itemize}
            \item \textbf{主修课程}：模式识别、计算机视觉、自然语言处理、数据结构、计算机组成原理、计算机网络\ 等。
            
        \end{itemize}
        
        \vspace{1.2em}
        \end{minipage}
    \end{minipage}


    % 实习经历
    \begin{minipage}[t]{\textwidth}
        \section[实习经历]{\makebox[\iconwidth][c]{\color{primary_color}{\faLaptopCode}}\quad 实习经历}
        \vspace{0.5em}

        {\large \textbf{字节跳动有限公司}} \hfill 2023 年 6 月 -- 2023 年 9 月 \\
        \textit{后端开发实习生} \hfill 北京
        \begin{itemize}
            \item \textbf{核心工作}：负责淘宝大促期间的流量网关优化，基于 Go 语言重构鉴权模块。
            \item \textbf{主要成果}：重构后的网关成功经受住双 11 流量洪峰的考验，通过优化锁粒度与内存管理，系统 延迟降低了 20\%，显著提升了用户体验。
            \item \textbf{技术栈}：Golang, Redis, Kubernetes, Prometheus。
        \end{itemize}
    
        \vspace{0.6em}
    
        {\large \textbf{腾讯科技有限公司}} \hfill 2022 年 12 月 -- 2023 年 3 月 \\
        \textit{\textbf{微信搜索团队}，算法实习生} \hfill 北京
        \begin{itemize}
            \item \textbf{核心工作}：参与 XX 大模型预训练数据处理，设计自动化清洗流水线。
            \item \textbf{主要成果}：清洗高质量数据超 100 万条，构建指令微调数据集，训练并持续提升模型视觉理解与通用推理
            能力，提升模型收敛速度 15\%。
            \item \textbf{技术栈}：Python, PyTorch, Spark, Hadoop。
        \end{itemize}
        
        \vspace{1.2em}
    \end{minipage}

    \begin{minipage}[t]{\textwidth}
    % 科研成果
    \section[科研成果]{\makebox[\iconwidth][c]{\color{primary_color}{\faAtom}}\quad 科研成果}

    % 科研著作（研究生）
    \textbf{This is One of Your Paper Published in Conference A.}
    \begin{itemize}
        \item \textbf{Mingzi Nide}, Daoshi Nide. \hfill 发表于 \textbf{ICLR}, \textbf{Oral}（CCF-A，人工智能顶级会议）
        \item 用一句话描述这份论文干了什么\dots\dots
    \end{itemize}
    \textbf{This is Another One of Your Paper Published in Journal B.}
    \begin{itemize}
        \item \textbf{Mingzi Nide}, Daoshi Nide. \hfill 发表于 \textbf{TPAMI}（一区TOP, JCR-Q1, IF11.2, CCF-A）
        \item 用一句话描述这份论文干了什么\dots\dots
    \end{itemize}
    \textbf{This is Another One of Your Paper Published in Journal C.}  \hfill 独立一作
    \begin{itemize}
        \item 发表于 \textbf{Information Fusion}（中科院一区, JCR-Q1, IF9.2, CCF-A） 
        \item 用一句话描述这份论文干了什么\dots\dots
    \end{itemize}
    \end{minipage}


    \newpage
    % 如有需要，可以添加额外的页面。不要忘记添加页眉页脚和背景相关的代码。
    % --- 页眉 ---
    \begin{tikzpicture}[remember picture, overlay]
        \node[anchor=north, inner sep=0pt](header) at (current page.north){
            \includegraphics[width=\paperwidth]{images/header.png}
        };
        \node[anchor=west](school_logo) at (header.west){
            \hspace{0.5cm}
            \includegraphics[width=0.25\textwidth]{images/banner.png}
        };
        \node[anchor=east](school_name) at(header.east){
            \textcolor{white}{\textbf{\school}}
            \hspace{0.5cm}
        };
    \end{tikzpicture}
    \vspace{-3.5em}

    % --- 页脚 ---
    \begin{tikzpicture}[remember picture, overlay]
        \node[anchor=south, inner sep=0pt](footer) at (current page.south){
            \includegraphics[width=\paperwidth]{images/footer.png}
        };
        % 联系方式
        \node[anchor=center] at(footer.center){\contact};
    \end{tikzpicture}
    
    % --- 背景 ---
    \begin{tikzpicture}[remember picture, overlay]
        \node[opacity=0.05] at(current page.center){
            \includegraphics[width=0.7\paperwidth, keepaspectratio]{images/logo.png}
        };
    \end{tikzpicture}

    % 3. 第二页的正文内容
    % 注意：因为 geometry 设置了 nohead 且我们手动控制了版心，这里通常不需要额外的 \vspace，
    % 但如果发现第二页顶部文字太靠上，可以适当加一点 \vspace*{...}
    
    % ================= 项目经历 =================
    \begin{minipage}[t]{\textwidth}
        % 图标建议：使用 faCode (代码) 或 faRocket (火箭/创新) 代表项目
        \section[项目经历]{\makebox[\iconwidth][c]{\color{primary_color}{\faCode}}\quad 项目经历}
        
        \vspace{0.5em}
        {\large \textbf{基于在线文档的考试系统}} \hfill 2023年9月--2023年12月 \\
        \textit{核心开发者} \hfill 独立项目
        \begin{itemize}
            \item \textbf{项目描述}：基于协同编辑技术，设计并实现支持多人实时在线作答与自动阅卷的考试系统。
            \item \textbf{主要工作}：利用 WebSocket 实现文档实时同步，采用 OT 算法解决冲突，结合 Redis 缓存提升阅卷效率。
            \item \textbf{成果}：支持2人以上同时进行跨学校考试无卡顿。
        \end{itemize}

        \vspace{0.6em}
        {\large \textbf{基于深度学习的图像识别算法}} \hfill 2023年3月--2023年6月 \\
        \textit{算法负责人} \hfill 实验室纵向项目
        \begin{itemize}
            \item \textbf{项目描述}：针对复杂光照条件下的目标检测任务，改进 YOLOv8 模型。
            \item \textbf{主要工作}：引入注意力机制模块，优化损失函数，提升小目标检测精度。
            \item \textbf{成果}：mAP 提升 3.5\%，相关成果已整理投稿至 CVPR 2024。
        \end{itemize}
        
        \vspace{1.2em}
    \end{minipage}

    % ================= 学生工作 =================
    \begin{minipage}[t]{\textwidth}
        % 图标建议：使用 faUsers (人群) 或 faUserTie (领导) 代表学生工作
        \section[学生工作]{\makebox[\iconwidth][c]{\color{primary_color}{\faUsers}}\quad 学生工作}
        
        \vspace{0.5em}
        {\large \textbf{校研究生会}}，主席 \hfill 2023年9月--至今
        \begin{itemize}
            \item \textbf{组织建设}：统筹管理校研会下设6个部门，定期召开部长例会，优化考核制度。
            \item \textbf{活动策划}：主导策划“学术文化节”、“十佳歌手大赛”等校级活动，累计参与人数超 2000 人。
        \end{itemize}

        \vspace{0.6em}
        {\large \textbf{北航对外国际合作部}}, 学生助理 \hfill 2022年9月--2023年6月
        \begin{itemize}
            \item \textbf{主要职责}：负责学院对外活动的节目编排与现场执行。
        \end{itemize}
        
        \vspace{1.2em}
    \end{minipage}
    
    
    \begin{minipage}[t]{\textwidth}
        \section[竞赛经历]{\makebox[\iconwidth][c]{\color{primary_color}{\faTrophy}}\quad 竞赛经历}
            \begin{itemize}
                \setlength{\itemsep}{0.5em} % 设置列表项间距
                \item \textbf{北航第33届“冯如杯”大赛，校级一等奖} \hfill 2020.09 -- 2021.09 \\
                      \textit{\textbf{职责}：团队负责人} %要写详细点就放这
                \item \textbf{国际大学生程序设计竞赛 (ICPC)，金牌 } \hfill 2023.09
            \end{itemize}
    \end{minipage}

    \vspace{1em}

    \begin{minipage}[t]{\textwidth}
        \section[所获荣誉]{\makebox[\iconwidth][c]{\color{primary_color}{\faStar}}\quad 所获荣誉}
        \vspace{0.5em}
        
        % 外层容器：使用 minipage 分为左右两栏
        \begin{minipage}[t]{0.48\linewidth}
            \begin{itemize}
                \setlength{\itemsep}{0.2em} % 压缩行间距，适应侧边栏
                \setlength{\leftmargin}{1.2em} % 调整缩进，防止文字太靠右
                \item \textbf{某年学业先进个人}
                \item \textbf{某大使} 
                \item \textbf{某竞赛一等奖}
            \end{itemize}
        \end{minipage}%
        % 注意：两个 minipage 之间不要有空行，否则会产生额外垂直间距
        \hfill % 自动填充中间空隙，使两列分开
        \begin{minipage}[t]{0.48\linewidth}
            \begin{itemize}
                \setlength{\itemsep}{0.2em}
                \setlength{\leftmargin}{1.2em}
                \item \textbf{某年某奖学金}
                \item \textbf{某年优秀团员}
                \item \textbf{某年某称号}
            \end{itemize}
        \end{minipage}
    \end{minipage}

    \vspace{1em}
    
    % 如果每行的内容不是很多，可以考虑使用 minipage，将内容分列展示
    \begin{minipage}[t]{\textwidth}
        \section[技能特长]{\makebox[\iconwidth][c]{\color{primary_color}{\faWrench}}\quad 技能特长}
        \begin{itemize}
        \setlength{\itemsep}{0.5em}
            \item 熟练使用 Python、Jvav、Rust 等编程语言。
            \item 熟练使用 Tensorflow、Pytorch 等深度学习框架。
        \end{itemize}
    \end{minipage}
    
    \vspace{1em}
    
    \begin{minipage}[t]{\textwidth}
        \section[自我评价]{\makebox[\iconwidth][c]{\color{primary_color}{\faComments}}\quad 自我评价}
        \begin{itemize}
            \setlength{\itemsep}{0.5em}
            \item \textbf{兴趣爱好：} 喜欢组乐队。
            \item \textbf{性格特点：} 精神状态良好。
        \end{itemize}
    \end{minipage}

\end{document}
